\chapter{مقدمه}\label{Chap:Chap1}
\minitoc
\setcounter{footnote}{0}
\pagestyle{plain}
در این فصل ابتدا به معرفی مساله‌ی یادگیری
\trans{بزرگ مقیاس، نمونه محدود}{Large-Sclae Few-Shot}
 پرداخته و ضمن معرفی دقیق، اهمیت و کاربردهای آن را مورد بررسی قرار می‌دهیم. در ادامه چالش‌های این حوزه را مورد بررسی قرار داده و بعد از معرفی هدف پژوهشی، ساختار این پایان‌نامه را شرح می‌دهیم.

\section{تعریف مسئله}
\label{sec:chap1:definition}
یادگیری عمیق در سالیان اخیر به طرز شگرفی توسعه پیدا کرده است و درهای جدیدی را در هوش مصنوعی گشوده است. امروزه بسیاری از مسائلی که تا قبل از معرفی یادگیری عمیق، چالش برانگیز بوده اند، دارای راه حل بدیهی هستند. در چند سال اخیر، رویکرد 
\trans{یادگیری مولد}{Generative Learning}
یک گام بزرگ دیگر در جهت پیشبرد اهداف هوش مصنوعی برداشته است. شاید بتوان گفت هوش مصنوعی و یادگیری عمیق در بین همه‌ی شاخه‌های علمی، به مرکز توجه پژوهشگران تبدیل شده است. از حدود دهه‌ی پنجاه میلادی که اولین
\trans{الگوی}{Paradigm}
هوش مصنوعی شکل گرفت و به الگوی 
\trans{سیستم‌های خبره}{Expert Systems}
، 
\trans{هوش مصنوعی نمادین}{Symbolic AI} 
یا همان تفکر منطقی و در نهایت الگوی 
\trans{یادگیری}{Learning}
که به صورت عملکرد منطقی تعریف می‌شود رسید، در تعریف هدف هوش مصنوعی دچار ابهام بوده ایم و به پدیده‌ای به نام 
\trans{اثر هوش مصنوعی}{AI Effect}
برخورد کرده ایم. به این معنی که همواره، رسیدن به یکسری اهداف مشخص به عنوان هوش مصنوعی تعریف شده است و وقتی به آن هدف رسیده ایم، گفته شده که آن موضوع دیگر هدف هوش مصنوعی نیست و یک هدف متعالی در نظر گرفته شده است. با این حال اکنون هدفی که برای هوش مصنوعی متصور هستیم رسیدن به 
\trans{هوش مصنوعی عمومی}{Artificial General Intelligence}
است که می‌توان آن را معادل هوش انسانی در نظر گرفت. در چند سال اخیر، محدودیت‌هایی که الگوریتم‌‌های یادگیری عمیق دارند بیشتر مورد توجه قرار گرفت که حاصل آن پدید آمدن رویکرد‌های یادگیری با الهام از یادگیری انسان است. یعنی علاوه بر اینکه در سطح عملکرد می‌خواهیم مشابه هوش انسانی عمل کنیم، در سطح الگوریتم یادگیری نیز فرآیند‌های یادگیری انسان را وارد کنیم. به این ترتیب می‌توانیم محدودیت‌هایی که در فرآیند آموزش یا استنتاج یک شبکه‌ی عمیق وجود دارند را تا حدی کاهش دهیم.

 از جمله‌ی این محدودیت‌ها 
\trans{داده ناکارآمدی}{Data Inefficiency}
 و 
 \trans{تعمیم‌پذیری}{Generalization}
 است. همواره یکی از مشکلات الگوریتم‌های یادگیری ژرف، بازدهی کم از لحاظ داده بوده است. یعنی برای آموزش یک مدل برای وظیفه‌ای خاص، مدل در صورتی به عملکرد قابل قبول می‌رسد که داده‌ی با حجم بالایی از آن مسئله را دیده باشد. این در حالی است که انسان در یادگیری همان 
 \trans{وظیفه}{Task}
 می‌تواند با مشاهده‌ی تعداد محدودی نمونه، مسئله را به خوبی یاد بگیرد. در نتیجه علی رغم تمامی پیشرفت‌های شگرف هوش مصنوعی، به خصوص از
 \trans{زمان پیدایش مجموعه داده‌ی بزرگ ایمیجنت}{ImageNet Moment}
  که افزایش توان محاسباتی به همراه عمیق شدن شبکه‌ها، روند هوش مصنوعی را به نقطه‌ی فعلی رساند، همواره دغدغه‌ی مجموعه‌ی دادگان آموزش برای یک مسئله مطرح بوده است.
  
  یک رویکردی که در ابتدا برای پاسخ‌گویی به مشکل داده ناکارآمدی شکل گرفت، گسترش مجموعه‌ی دادگان به طرق مختلف بود. یعنی در ابتدا به عنوان یک اصل پذیرفته شد که مدل‌های یادگیری ژرف 
  \trans{پیچیدگی نمونه‌ای}{Data Complexity}
  بالایی دارند، اما سعی شد از روی مجموعه‌ی دادگان برچسب‌خورده‌ی فعلی، داده‌های بیشتری تولید و به مدل برای آموزش داده شود
  \cite{Shorten2019ASO}
  .
  به این دسته روش‌های مبتنی بر 
  \trans{داده‌افزایی}{Data Augmentation}
  می‌گویند. نکته‌ی قابل توجه این است که در عصر امروزه، برای اکثر مسائل، داده‌ی خام به تعداد زیاد در دسترس است اما تمیز کردن این داده و برچسب‌زدن آن برای آموزش یک شبکه در یک مسئله‌ی خاص مشکل و هزینه‌بر است. در طول سالیان گذشته، رویکردهای دیگری برای حل این مشکل توسعه داده شده است. یک رویکرد استفاده از شبکه‌های مولد است
  \cite{Wang2018LowShotLF,Hariharan2016LowShotVR}
  .
   در شبکه‌های مولد سعی میشود با یادگیری
   \trans{توزیع توأمان}{Joint Distribution}
    داده‌ها و برچسب‌ها، قدرت تولید نمونه‌ی جدید را بدست آورد. در نتیجه از آن می‌توان به
  عنوان یک راه مقابله با کمبود داده بهره برد. راه حل دیگر جدایی از یادگیری نظارت شده است. در این زمینه روشهایی مثل یادگیری
  \trans{نیمه‌نظارتی}{Semi-Supervised}
  و یادگیری
  \trans{خودنظارتی}{Self-Supervised}
   مطرح شده اند
   \cite{Jing2019SelfSupervisedVF,Liu2020SelfSupervisedLG,Reddy2018SemisupervisedLA,Kingma2014SemisupervisedLW}
   .
   
   در یادگیری نیمه‌نظارتی سعی می‌شود از ترکیب یادگیری نظارت‌شده و یادگیری بدون
   ناظر استفاده شود. به این صورت که بخش کوچکی از داده برچسب‌گذاری شود و در کنار بخش بزرگتری از داده‌ی برچسب نخورده، در آموزش مدل  استفاده شود. یک روش مشهور در این حوزه
   \trans{خوشه‌بندی}{Clustering}
    داده‌ها در فضا و تعیین برچسب هر خوشه با توجه به نزدیکی به داده‌های برچسب‌دار است
    \cite{Bair2013SemisupervisedCM,Qin2019ResearchPO}
    .
   در یادگیری خودنظارتی، مدل سعی می‌کند یک علامت نظارتی از درون ساختار داده بدست آورد. یک روش مطرح در این حوزه
   \trans{یادگیری تباینی}{Contrastive Learning}
    است. در این نوع یادگیری، با بهره‌بردن از داده‌افزایی، نمونه‌های مثبت از روی یک داده تولید می‌شود. سایر نمونه‌های موجود در
    \trans{دسته}{Batch}
     به عنوان نمونه‌ی منفی در نظر گرفته می‌شوند و
     \trans{تابع خطا}{Loss Function}
      برای مدل به نوعی تعریف می‌شود که فاصله‌ی بین نمونه‌های مثبت را کمینه کند و همزمان فاصله‌ی بین یک نمونه‌ی مثبت و منفی بیشینه شود
    \cite{Chen2020BigSM,Chen2020ASF}
    .
    همه‌ی این تلاش‌ها برای مقابله با کمبود داده‌های برچسب‌دار یک نقطه‌ی ضعف مشترک دارند. وقتی تعداد اندکی داده در اختیار باشد، گسترش این دادگان به هر نحوی موجب به وجود آمدن
    \trans{سوگیری}{Bias}
     است. چرا که یک نمونه‌ی کوچک از یک توزیع توانایی بازنمایی بخش کوچکی از آن توزیع را دارد و وقتی برپایه‌ی آن دادگان جدید تولید شود، سوگیری به سمت آن بخش از توزیع به وجود می‌آید
     \cite{Yang2021FreeLF}
     .
  
  از طرف دیگر، شبکه‌های عمیق یک سیستم خبره هستند و روی حل یک مسئله‌ی خاص تمرکز دارند. این موضوع نه تنها با هدف هوش مصنوعی عمومی در تناقض است، بلکه با مسائل دنیای واقعی تطابق کامل ندارد. چرا که برای حل یک مسئله‌ی حقیقی، باید آن را به بخش‌های کوچکتر خرد کرد و برای هر کدام از این بخش‌ها با فراهم کردن مجموعه دادگان برچسب‌گذاری شده، برای آن وظیفه‌ی خاص یک مدل خبره آموزش داد. رویکرد آموزش شبکه‌های ژرف به صورت احتمالاتی است و در نگاه 
  \trans{تمایزگذار}{Discriminative}
  همواره سعی میکند احتمال یک برچسب به شرط داده‌ی فعلی را تخمین بزند. بنابراین با فرض وجود دادگان کافی، وقتی تعداد دسته‌ها را زیاد میکنیم، احتمال‌های
  تخصیص داده شده به هر دسته با فاصله‌ی کمی از هم قرار می‌گیرند و با کاهش
  \trans{مرز تصمیم‌گیری}{Decision Boundary}
  برای مدل، خطای آن بالا می‌رود. در نتیجه با  بالا رفتن تعداد دسته‌ها، دقت مدل‌ها به طور محسوسی افت میکند. دسته‌بندی بزرگ مقیاس همچنان به عنوان یکی از چالش ‌های مهم این حوزه مطرح است. 
  
  
  \begin{table}[h]
	\caption{دسته‌بندی مسائل یادگیری و تمرکز پژوهش فعلی}
	\label{table:problemTypes}
	\centering
	%  	\resizebox{\columnwidth}{!}{%
		\begin{tabular}{@{}ccc@{}}
			\toprule
			& تعداد کلاس کم           &تعداد کلاس زیاد             \\
			 \midrule
			نمونه محدود            & روش‌های \lr{Few-Shot}                 
			&\textbf{تمرکز این پژوهش} 
			\\
			نمونه زیاد & \multicolumn{2}{c}{تمرکز مدل‌های متداول یادگیری ژرف}  \\ \bottomrule
		\end{tabular}%
		%  	}
\end{table}

اگر برای مسائل یک افراز مطابق جدول 
\ref{table:problemTypes}
در نظر بگیریم، تمرکز این پژوهش روی مسائلی با تعداد کلاس زیاد و بزرگ مقیاس به همراه تعداد کم نمونه‌ی آموزش برای هر کلاس است. برای مسائل نمونه محدود، در چند سال اخیر راهکار‌های مختلفی پیشنهاد داده شده است که یکی از مهم‌ترین آن‌ها رویکرد 
\trans{متایادگیری}{Meta-Learning}
است. در این رویکرد سعی می‌شود یادگیری برای فرآیند یادگیری اتفاق بیفتد و به نوعی تمرکز روی استفاده از تجربیات مشترک در مسائل یادگیری از قبل دیده شده است. در این دسته از راهکار‌ها، معمولا تعداد کلاس‌های هر مسئله‌ی یادگیری را کمتر از ۲۰ در نظر می‌گیرند، چرا که این روش‌ها در مسائل با تعداد دسته‌ی بالا نمی‌توانند به عملکرد مطلوب برسند. از طرف دیگر، در مسائلی که نمونه زیاد آموزش وجود دارد، رویکرد‌های مرسوم یادگیری ژرف قابل استفاده است. اما در این حوزه هم توجه به مسائلی مانند خاصیت ذاتی 
\trans{مجموعه باز}{Open-Set}
بودن دسته‌بندی بزرگ مقیاس حائز اهمیت است. چرا که وقتی با یک مسئله‌ی دسته‌بندی چند هزار دسته‌ای روبرو هستیم، احتمال بالایی وجود دارد که به صورت مستمر به دسته‌های موجود اضافه شود. بنابراین یکی از چالش‌های اصلی این دسته از مسائل، بحث تعمیم‌پذیری و داشتن نگاه مولد به جای نگاه تمایزگذار است. 

در مسائل بزرگ مقیاس، به دلیل تعدد دسته‌ها، می‌توان آن‌ها را در قالب یک
\trans{سلسله‌مراتب}{Hierarchy}
 از 
\trans{درشت دانه}{Coarse Grained}
به 
\trans{ریزدانه}{Fine Grained}
مرتب کرد و از اطلاعات موجود در سلسله‌مراتب ذاتی دسته‌ها استفاده کرد. همان‌طور که انسان دنیا را به صورت سلسله‌مراتبی درک می‌کنند
\cite{Hubel1968ReceptiveFA,Fukushima1980NeocognitronAS,Kuzovkin2017ActivationsOD}
، چنین خاصیت ذاتی سلسله‌مراتبی در محتوای بصری مورد بررسی قرار گرفته‌است تا در طراحی معماری مدل‌های مختلف کمک کند. 
بنابراین، استدلال ما این بوده که در تنظیمات پرچالشی مثل دسته‌بندی بزرگ مقیاس، نمونه محدود، که نیازمند درک و استدلال ظریف است، با کمبود داده، نیاز حیاتی به استفاده از ماهیت سلسله‌مراتبی محتوای بصری است تا از اطلاعات ضمنی موجود بتوان بیشترین استفاده را برد.

در پژوهش پیش رو، تمرکز اصلی روی حل کردن یک مسئله‌ی دسته‌بندی بزرگ مقیاس با تعداد نمونه‌ی محدود آموزش برای هر دسته با نگاه سلسله مراتبی است.
به صورت خلاصه، مهمترین دستاوردهای پژوهش ما را می‌توان در موارد زیر خلاصه نمود:

\begin{itemize}
\item 
ما یک روش یادگیری نوین برای استفاده از سلسله‌مراتب ذاتی موجود در دسته‌ها با طراحی یک
\trans{سوگیری القایی}{Inductive Bias}
مبتنی بر آن،‌ برای حل مسئله‌ی یادگیری بزرگ مقیاس، نمونه محدود پیشنهاد داده‌ایم.
\item
از مدل‌های 
\trans{بنیادی}{Foundation}
\trans{پیش‌آموزش دیده}{Pre-Trained}
برای بدست‌آوردن 
\trans{بازنمایی}{Representation}
غنی برای رسیدن به داده کارآمدی و افزایش مقاومت در برابر 
\trans{تغییر توزیع}{Distribution Shift}
استفاده کرده‌ایم.

\item
یک رویکرد کاربردی برای کاهش 
\trans{تغییر زیر جمعیت}{Sub-Population Shift}
در سطح کلاس‌های درشت‌دانه یا همان 
\trans{ابرکلاس}{Superclass}
ها ارائه داده‌ایم.
\item
مدل معرفی شده در این پژوهش حاوی المان‌های سبک و کارآمد است که هم امکان پیاده‌سازی روی سخت‌افزار محدود در شرایط خاص مسئله یادگیری بزرگ مقیاس را فراهم می‌کند و هم نسبت به روش‌های پیشین از 
\trans{بهینگی محاسباتی}{Computational Efficiency}
 بهتری بهره می‌برد.
\item
ما از یک نگاه مولد استفاده کرده‌ایم و در قالب یک چارچوب 
\trans{بیزی}{Bayesian}
به ترکیب اطلاعات حاصل از سطوح مختلف سلسله‌مراتب می‌پردازیم.
\item 
ما نشان داده‌ایم که راهکار پیشنهادی می‌تواند به سادگی در تنظیمات 
\trans{بدون یادگیری}{Free-Training}
 به کار گرفته شود.
\end{itemize}

\section{رویکرد حل مسئله}\label{sec:approach}

مسئله‌ی تعریف شده را در قالب یک یادگیری سلسله‌مراتبی در دو سطح مجزا تعریف می‌کنیم. در سطح اول که سطح درشت‌دانه می‌باشد، ابردسته‌ها قرار دارند و در سطح دوم هم کلاس‌های ریزدانه که به آن‌ها 
\trans{دسته پایه}{Base Class}
می‌گوییم قرار می‌گیرند. برای هر سطح یک دسته‌بند مجزا آموزش داده می‌شود و در استنتاج دسته‌بند سطح دوم، از نتایج دسته‌بند سطح اول استفاده ‌می‌شود. در واقع کل مسئله یادگیری سلسله‌مراتبی را در قالب یک مسئله‌ی دسته‌بندی 
\trans{یکی از همه برنده}{Winner Takes All}
تعریف می‌کنیم. در سطح اول، در دامنه‌ی مدنظر، ابرکلاس‌هایی قرار می‌گیرند که هیچ‌گاه جمعیت آن‌ها تغییری نخواهد کرد. به عنوان مثال در یک مسئله‌ی دسته‌بندی گونه‌های جانوری، با یک دسته‌بندی بزرگ مقیاس مواجه هستیم. در سطح اول دسته‌هایی مانند پرندگان و خزندگان و ... قرار می‌گیرند. در حالی که در سطح دوم ریزدانه‌ترین گونه‌هایی که برای حیوانات متصور هستیم قرار می‌گیرند. با توجه به نمونه محدود بودن مسئله، وقتی آن را در این قالب بازتعریف می‌کنیم، برای دسته‌های سطح اول به تعداد قابل توجهی نمونه آموزش تجمیع می‌شود. چرا که نمونه‌های آموزشی هر ابردسته به اندازه‌ی دسته‌های پایه فرزند آن ضربدر تعداد نمونه‌های آموزشی هر کدام از آن کلاس‌های ریزدانه است. 

\subsection{معماری مدل سطح درشت‌دانه}
برای سطح اول، یک دسته‌بند تمایزگذار در نظر گرفته‌ایم که علاوه بر آموزش نظارت شده، به صورت نیمه نظارتی هم 
\trans{تنظیم بهینه}{Fine-Tune}
می‌شود. علت این انتخاب ثابت بودن ابردسته‌ها در یک دامنه است. در این صورت دسته‌ای به آن‌ها به صورت مستمر اضافه نمی‌شود و به همین دلیل نیاز به باز آموزش یک شبکه‌ی تمایزگذار یا استفاده از رویکرد مولد نیست. از طرفی، با تجمیع نمونه‌های آموزشی دسته‌های پایه و عملکرد بهتر دسته‌بند تمایزگذار نسبت به مولد به صورت کلی، این انتخاب صورت گرفته است. 

\subsection{معماری مدل سطح ریز‌دانه}
برای سطح دوم، از یک رویکرد متایادگیری به نام شبکه‌های پروتوتیپیکال
\cite{Snell2017PrototypicalNF}
 استفاده کرده‌ایم. این شبکه یک دسته‌بند مولد است که می‌تواند مشکلات ناشی از اضافه شدن کلاس‌های ریزدانه و مجموعه‌باز بودن را حل کند. از طرفی از خانواده‌ی الگوریتم‌های متایادگیری است که که برای تنظیمات نمونه محدود مناسب است و دارای ویژگی‌های مانند تطبیق سریع، تعمیم‌پذیری، کشف ساختار‌های مشترک و اساسی است. شبکه‌ی پروتوتیپیکال یک دسته‌بند 
\trans{ناپارامتری}{Non-Parametric}
 است که مرحله‌ی 
 \trans{یادگیری سریع}{Fast Learning}
 در الگوریتم‌های متایادگیری را بدون آپدیت خاصی و تنها با محاسبه‌ی پروتوتایپ‌ها بدست می‌آورد. به همین دلیل از لحاظ محاسباتی کارآمد است و نیازی به نگه داشتن 
 \trans{گراف محاسباتی}{Computational Graph}
 سنگین و محاسبه‌ی مشتقات مراتب بالاتر و ماتریس 
 \trans{هسین}{Hessian}
 ندارد. در سطح دوم، شبکه‌ی پروتوتیپیکال را به صورت 
 \trans{مشروط}{Conditioned}
 روی پیش‌بینی شبکه‌ی سطح اول آموزش می‌دهیم و به صورت مشروط از آن استنتاج می‌گیریم.
 
 \subsection{تنظیمات متاآموزش و متاتست}
 
 تنظیمات آموزش و استنتاج را به صورت مراحل 
 \trans{متاآموزش}{Meta Train}
 و 
 \trans{متاتست}{Meta Test}
 می‌شکنیم. در مسائل بزرگ مقیاس، تغییر زیر جمعیت متداول است. یعنی همانطور که گفته شد، در سطح ابردسته‌ها کلاس‌ها ثابت می‌مانند اما در سطح دسته‌های پایه، دسته‌ها کاملا
 \trans{بدیع}{Novel}
  و 
  \trans{مجزا}{Disjoint}
   از دسته‌های پایه اولیه می‌شوند. این تنظیمات برای بدست‌آمدن 
 \trans{مقاومت}{Robustness}
 نسبت به مجموعه‌باز بودن و پیدایش دسته‌های پایه جدید در نظر گرفته شده‌است. به این ترتیب، در دو فاز متاآموزش و متاتست، در سطح ابردسته‌ها، دسته‌هایی کاملا یکسان وجود دارد اما در سطح ریزدانه و دسته‌های پایه، دسته‌ها کاملا مجزا و متفاوت از هم هستند.
 
 برای آموزش سطح اول، با توجه به داده‌های موجود در فاز متاآموزش، یک دسته‌بند به صورت یادگیری با ناظر آموزش داده می‌شود. در واقع در مرحله‌ی متاآموزش، باید برای هر داده برچسب ریزدانه و درشت‌دانه‌ی آن را بدانیم تا بتوانیم دسته‌بند ابرکلاس‌ها را به صورت با ناظر آموزش دهیم. اما در فاز متاتست، برای داده‌های کلاس‌های بدیع، هیچ اجباری روی دانستن برچسب درشت‌دانه داده‌ها نداریم و از این برچسب هیچ استفاده‌ای نمی‌کنیم. اما با توجه به اینکه تعداد محدود نمونه‌ی 
 \trans{پشتیبان}{Support}
 هر کلاس پایه را در اختیار داریم، از آن‌ها برای تنظیم بهینه دسته‌بند سطح اول می‌توانیم استفاده کنیم. چرا که به صورت ضمنی وقتی یک تعدادی داده متعلق به یک کلاس ریزدانه هستند، حتما متعلق به یک کلاس درشت‌دانه هم هستند اما برچسب اصلی ابرکلاس آن‌ها را نمی‌دانیم. به همین دلیل با یک روش نیمه نظارتی به تنظیم بهینه وزن‌های دسته‌بند سطح اول در مرحله‌ی متاتست می‌پردازیم.
 
  نکته‌ی حائز اهمیت این است که در معماری دسته‌بند‌های هر دو سطح، از یک 
 \trans{شبکه‌ی پایه}{Backbone}
 پیش آموزش دیده از مدل‌های بنیادی استفاده می‌کنیم به که استفاده از دانش و بازنمایی 
 \trans{بیانگر}{Expressive}
 و داده کارآمدی مناسب در عین توان محاسباتی پایین در مرحله‌ی آموزش شبکه‌ها برسیم.
 
 در آموزش دسته‌بند سطح دوم، در فاز متاآموزش، یک شبکه‌ی 
 \trans{کاملا متصل}{Fully Connected}
 را به صورت ایپزودیک و با ساختن وظیفه‌های متعدد مشابه روش پروتوتیپیکال آموزش می‌دهیم. ورودی به شبکه، حاصل پیش‌بینی احتمال تعلق داده‌ی فعلی به هر یک از ابردسته‌ها می‌باشد. یعنی در هنگام آموزش، شبکه‌ی سطح دوم روی پیش‌بینی سطح اول مشروط می‌شود. در فاز متاتست، در یک وظیفه که به اندازه‌ی دسته‌های ریزدانه 
 \trans{راه}{Way}
 دارد، از این سطح به صورت مشروط بر تخمین سطح اول استنتاج می‌گیریم. خروجی این سطح رویکرد مولد دارد و در واقع به شرط هر کدام از کلاس‌های ریزدانه، احتمال وقوع داده‌ی فعلی را خروجی می‌دهد. به همین دلیل است که بدون آموزش مجدد در فاز متاتست و با وجود پیدایش کلاس‌های کاملا بدیع می‌توانیم از آن استنتاج بگیریم.
 
 در نهایت در یک ستینک بدون یادگیری هم آزمایش‌هایی انجام داده‌ایم که نشان داده شده است از این چارچوب ارائه شده می‌توان بدون هیچ آموزشی برای دسته‌بند سطح دوم در فاز متاآموزش و تنها با استفاده از یک شبکه‌ی پایه پیش‌آموزش داده شده هم استفاده کرد.
 
\section{چالش‌ها}\label{sec:challenges}

در این بخش چالش‌های اصلی که الگوریتم‌های یادگیری مستمر با آن‌ها درگیر هستند را معرفی می‌کنیم.

\subsection{محدودیت تعداد کلاس  در روش‌های نمونه محدود}
در روش‌های نمونه محدود متداول، مثل روش‌ 
MAML
\cite{Finn2017ModelAgnosticMF}
 یا 
 SNAIL
 \cite{Mishra2017ASN}
  بر روی تعداد دسته‌های هر وظیفه محدودیت وجود دارد. در این روش‌ها، حداکثر دسته‌هایی که برای یک وظیفه در نظر می‌گیرند حداکثر حدود ۲۰ کلاس است. بنابراین در مسئله‌ی تعریف شده که نیاز به دسته‌بندی نمونه محدود بین چند هزار دسته هست، نمی‌توان از این روش‌ها به صورت مستقیم استفاده کرد.
  
\subsection{مجموعه‌باز بودن دسته‌بندی بزرگ مقیاس}

وقتی مسئله دارای تعداد زیادی کلاس است، به صورت ذاتی قابلیت اضافه شدن دسته‌های جدید وجود دارد. از طرفی یکی از تنظیمات متداول در مسائل بزرگ مقیاس، تغییر زیر توزیع است. این موضوع با تغییر توزیع متداول کاملا متفاوت است. در تغییر توزیع، مثلا از حیطه‌ی تصاویر طبیعی به تصاویر طراحی یا نقاشی شده می‌رویم. اما در تغییر زیر توزیع، در همان توزیع هدف، با یک سری دسته‌ی دیگر روبرو می‌شویم که دقیقا از همان توزیع می‌آیند و تغییرات کمی دارند. بنابراین در اینگونه مسائل باید یک رویکرد مولد در نظر گرفت.

\subsection{انتشار خطا در سلسله‌مراتب}
وقتی از سلسله‌مراتب استفاده می‌‌کنیم، در پیدا کردن تعداد سطوح بهینه دچار یک 
\trans{مصالحه}{ُTrade Off}
هستیم. چرا که با بالارفتن سطوح سلسله‌مراتب، خطای یک سطح به تمامی سطوح پایینی انتشار پیدا می‌کند. یعنی در واقع در مسیر رسیدن به یک برگ، هر کدام از دسته‌بند های سطوح بالایی اشتباه کنند، پیش‌بینی دسته‌بند‌های دیگر هر چه باشد غلط است چون مسیر استنتاج در مسیر غلطی می‌افتد. این مشکل وقتی نمود پیدا می‌کند که در هنگام استفاده از روش یکی از همه برنده، دسته‌بند به ازای هر گره داشته باشیم. برای حل این مشکل راهکاری که ما اتخاذ کردیم، استفاده از دسته‌بند در هر سطح است. یعنی به جای اینکه نگاه کنیم اکنون در کدام گره هستیم و فقط بین فرزندان آن گره دسته‌بندی انجام دهیم، به ازای هر سطح، یک دسته‌بند داریم که بین گره‌های موجود در آن سطح دسته‌بندی می‌کند. حال برای ترکیب اطلاعات حاصل از دسته‌بندی سطح قبلی، با یک دیدگاه بیزی به ترکیب این اطلاعات می‌پردازیم که در واقع مشروط شدن تصمیم‌گیری دسته‌بند هر سطح براساس خروجی دسته‌بند سطح قبلی است. با این اوصاف، به این نتیجه رسیدیم که تعداد بهینه سطوح برای رسیدن به بالاترین دقت، ۲ سطح است. این انتخاب در کنار در نظر گرفتن جوانب دیگر مثل ثابت بودن ابردسته‌ها در عین داشتن تغییر زیرجمعیت بوده است. یعنی برای سطوح بالایی که می‌خواهیم دیتا به اندازه‌ی کافی تجمیع شود و به صورت تمایزگذار آموزش داده شود، یک سطح ابردسته کافی است. 

از طرف دیگر، تمییز دادن کلاس‌های مختلف در فضای ویژگی آموخته شده با صدها دسته با نمونه‌های اندک، ما را به سمتی برد که از سلسله‌مراتب استفاده کنیم. چرا که استفاده از سلسله‌مراتب، تعداد بالای دسته را به سمت مسائل دسته‌بندی متوازن‌تر با مرتبه‌ی لگاریتمی می‌برد.


\section{اهمیت و کاربرد}
در حال حاضر، قالب رایج در آموزش شبکه‌های عصبی برای استفاده‌های کاربردی به این صورت است که ابتدا یک دیتاست نسبتاً مفصل متناسب با وظیفه مد نظر جمع آوری می‌شود، سپس مدل با این داده‌ها آموزش می‌بیند و پس از منجمد کردن وزن‌ها، مدل آماده به کارگیری در محیط خواهد بود. این رویکرد معایب جدی دارد: اولاً‌ استفاده از این رویکرد قابلیت مقیاس پذیری را پایین می‌آورد چرا که جمع‌آوری مجموعه دادگان غنی و به خصوص نگه‌داری آن‌ها می‌تواند به شدت پر هزینه و زمان‌گیر باشد. ثانیاً قابلیت تطبیق‌پذیری با شرایط محیطی در این رویکرد به حداقل می‌رسد چرا که مدل بعداً در صورت مواجهه با تجارب جدید، اجازه تغییر در پارامتر‌های خود برای تطبیق‌پذیری بهتر با محیط را ندارد. از سویی دیگر امکان آموزش مجدد مدل با استفاده از کل دادگان وجود ندارد چرا که این امر بسیار زمان‌گیر خواهد بود
\cite{parisi}.
 لذا مجهز کردن عامل‌های هوشمند به توانایی یادگیری مستمر، یک گام بسیار اساسی در جهت عمومیت بخشی و پدید آوردن خودمختاری در این عامل‌هاست. برای پی بردن به اهمیت بحث یادگیری مستمر بیشتر، می‌توان ردپای آن را در کاربردهای مختلف دنبال کرد. 
\begin{itemize}
		\item 
		\textbf{رباتیک و عامل‌های هوشمند:}
با پیشرفت یادگیری تقویتی در سالیان اخیر، تحقیق و توسعه ربات‌های هوشمند اهمیت بسیار زیادی پیدا کرده است. ربات‌ها باید بتوانند از محیط پیرامون خود یاد بگیرند و همزمان از طریق کنش با محیط و انجام مشاهدات پیوسته، رفتار خود را بهبود دهند \cite{robocl}.
از این رو
\trans{یادگیری تقویتی مستمر}{Continual Reinforcement Learning}
از جمله مباحث پژوهشی نوینی است که اخیراً‌ مورد توجه بیشتری قرار گرفته است
 \cite{ContinualRL,CRL}.

\item 

\textbf{سیستم‌های توصیه کننده:}
\trans{سیستم‌های توصیه کننده }{Recommender Systems}
سیستم‌هایی هستند که هدف آن‌ها معرفی محصولات، متن‌ها، تبلیغات یا هر مطلب دیگر متناسب با سلیقه هر کاربر است. در این سیستم‌ها شخصی‌سازی متناسب با سلایق کاربر، از طریق کنش‌ها و تعاملات تدریجی الگوریتم با هر کاربر اتفاق می‌افتد. لذا تطبیق سریع مدل با علایق کاربر اهمیت خیلی زیادی دارد. از سوی دیگر آموزش مدل روی کل مجموعه دادگان به ازای هر کاربر جدید امری نشدنی است. لذا ویژگی تطبیق پذیری بهینه از طریق تعامل با محیط یک مشخصه اجتناب ناپذیر برای این سیستم‌هاست \cite{Recom}.
		
		\item 
		\textbf{ سیستم‌های نهفته:}
		امروزه ابزارها و تراشه‌های الکترونیکی از جمله تلفن‌های همراه در میان مردم رواج زیادی پیدا کرده‌اند. لذا بسیاری از محققین در زمینه پیاده‌سازی بهینه الگوریتم‌های یادگیری ماشین بر روی تراشه‌ها و 
		\trans{سیستم های نهفته}{Embedded Systems}
		 فعالیت می‌کنند. از سویی به دلیل محدودیت‌های سخت‌افزاری از جمله محدودیت فضای ذخیره‌سازی‌ داده در کنار چیپ، بسیاری از الگوریتم‌های یادگیری ماشین با مشکلاتی مواجه می‌شوند. در الگوریتم‌های یادگیری مستمر فرض بر این است که فضای ذخیره‌سازی اضافه‌ای وجود ندارد یا این که این فضا بسیار کوچک است. لذا این الگوریتم‌ها کاملاً با محدودیت‌های عملی پیاده‌سازی میکروچیپ‌ها سازگار خواهند بود\cite{embCL}. 


\end{itemize}

\section{هدف پژوهش}


با توجه به مباحث ذکر شده، در این پژوهش ما به دنبال ارائه یک روش یادگیری مستمر هستیم که در عین صفر بودن فراموشی و استفاده حداقلی از حافظه جانبی، اجازه تطبیق‌پذیری به کلاس‌ها و داده‌های جدید را فراهم نماید. برای این منظور، از بستر الگوریتم‌های 
\trans{یادگیری متا-مستمر}{Meta-Continual Learning}
استفاده کرده و با الهام از یادگیری انسان در حوزه علوم اعصاب شناختی، یک الگوریتم کارآمد در این بستر معرفی می‌کنیم.

\section{ساختار پایان‌نامه}
در ادامه‌ی مستند حاضر، در فصل دوم به تشریح کار‌های پیشین در زمینه‌ی یادگیری مستمر، 
\trans{متا-یادگیری}{Meta-Learning}
 و یادگیری متا-مستمر پرداخته و ضمن آوردن شواهد زیستی متناسب با این روش‌ها، آن‌ها را از جنبه‌های مختلف با یکدیگر مقایسه می‌کنیم.
در فصل سوم، راهکار خود را که الهام گرفته از نحوه یادگیری در مغز انسان است را برای بهبود یادگیری متا-مستمر ارائه می‌کنیم و در فصل چهارم، با ارائه‌ آزمایش‌های متعدد موثر بودن این روش‌ را نمایش می‌دهیم. در نهایت ایرادها، چالش‌ها و مسیر پیش‌رو برای ادامه تحقیقات را در فصل آخر این پایان‌نامه ارائه می‌کنیم.



